\section*{TD 6 --- SVM et noyaux}
\addcontentsline{toc}{section}{TD 6 --- SVM et noyaux}

\subsection*{Exercice 1 --- Marge et vecteurs de support \hfill \normalfont\footnotesize (8 min)}
Deux classes sont séparables par une droite. Trois points sont sur la marge :
$A$, $B$ et $C$.
\begin{enumerate}
  \item Qu'est-ce qu'un vecteur de support ?
  \item Que devient la frontière si l'on supprime un point très éloigné de la
        marge ?
  \item Et si l'on supprime $A$ ?
\end{enumerate}

\subsection*{Exercice 2 --- Effet de $C$ \hfill \normalfont\footnotesize (8 min)}
On fait varier $C$ de $0{,}01$ à $100$.
\begin{enumerate}
  \item Que se passe-t-il pour $C$ très petit ?
  \item Et pour $C$ très grand ?
  \item Comment choisir $C$ ?
\end{enumerate}

\subsection*{Exercice 3 --- Noyau RBF \hfill \normalfont\footnotesize (10 min)}
Deux points : $x_1=(1,2)$ et $x_2=(3,4)$. Noyau
$K(x_1,x_2)=\exp(-\gamma\lVert x_1-x_2\rVert^2)$.
\begin{enumerate}
  \item Calculer $\lVert x_1-x_2\rVert^2$.
  \item Calculer $K$ pour $\gamma=0{,}1$, puis $\gamma=0{,}5$, puis $\gamma=2$.
  \item Que se passe-t-il quand $\gamma$ augmente ?
\end{enumerate}

\subsection*{Exercice 4 --- SVM ou arbre ? \hfill \normalfont\footnotesize (8 min)}
Indiquer le modèle le plus adapté et pourquoi :
\begin{enumerate}
  \item 80 observations, 500 variables, frontière nette.
  \item 200\,000 lignes, variables catégorielles et manquantes.
  \item Besoin d'expliquer chaque décision à un client.
\end{enumerate}

\subsection*{Exercice 5 --- Que fait le noyau ? \hfill \normalfont\footnotesize (6 min)}
\begin{enumerate}
  \item Pourquoi ne calcule-t-on pas explicitement la transformation
        $\varphi(x)$ dans un espace de grande dimension ?
  \item Que permet l'astuce du noyau ?
  \item Citer deux noyaux usuels.
\end{enumerate}


\subsection*{Exercice 6 --- Calculer une marge \hfill \normalfont\footnotesize (8 min)}
Le séparateur est $w=(1,1)$, $b=0$.
\begin{enumerate}
  \item Calculer $\lVert w\rVert$.
  \item Calculer la marge $2/\lVert w\rVert$.
  \item Calculer la distance du point $(1,1)$ à la droite. Que remarque-t-on ?
\end{enumerate}

\subsection*{Exercice 7 --- Points mal classés et marge souple \hfill \normalfont\footnotesize (7 min)}
Deux points sont du mauvais côté de la frontière.
\begin{enumerate}
  \item Que se passe-t-il si l'on exige une séparation parfaite ?
  \item Que permet la variable ressort (\emph{slack}) ?
  \item Quel paramètre règle le compromis marge / erreurs ?
\end{enumerate}

\subsection*{Exercice 8 --- Noyau polynomial \hfill \normalfont\footnotesize (6 min)}
Noyau $K(x,x')=(1+x\cdot x')^2$, avec $x=(1,2)$ et $x'=(2,1)$.
\begin{enumerate}
  \item Calculer $x\cdot x'$.
  \item En déduire $K(x,x')$.
  \item Pourquoi ne calcule-t-on pas explicitement $\varphi(x)$ ?
\end{enumerate}

\subsection*{Exercice 9 --- Complexité \hfill \normalfont\footnotesize (6 min)}
Un SVM exact a un coût de l'ordre de $O(n^2)$ à $O(n^3)$.
\begin{enumerate}
  \item Que cela implique-t-il avec 500\,000 observations ?
  \item Citer deux alternatives pratiques.
  \item Le noyau RBF change-t-il cet ordre de grandeur ?
\end{enumerate}

\newpage
